\documentclass{beamer}
\usepackage{bm}
\usepackage{tabularx}
\setbeamertemplate{caption}[numbered]
\usepackage{xcolor}
\usepackage{multirow}
\usepackage{movie15}
\usepackage{Sweave}
%\usepackage{nameref}% http://ctan.org/pkg/nameref
%\newcommand{\currentchapter}{}
%\let\oldchapter\chapter
%\renewcommand{\chapter}[1]{\oldchapter{#1}\renewcommand{\currentchapter}{#1}}
\renewcommand{\thesection}{4.\arabic{section}}
\newcommand{\currentsection}{}
\let\oldsection\section
\renewcommand{\section}[1]{\oldsection{#1}\renewcommand{\currentsection}{#1}}

\newcommand{\currentsubsection}{}
\let\oldsubsection\subsection
\renewcommand{\subsection}[1]{\oldsubsection{#1}\renewcommand{\currentsubsection}{#1}}

\newcommand{\currentsubsubsection}{}
\let\oldsubsubsection\subsubsection
\renewcommand{\subsubsection}[1]{\oldsubsubsection{#1}\renewcommand{\currentsubsubsection}{#1}}


\newenvironment{wideitemize}%
  {\begin{itemize}%
%    \setlength{\itemsep}{0pt}%
    \setlength{\parskip}{1em}}%
  {\end{itemize}}
%\usetheme{CambridgeUS}
\usetheme{Madrid}

%\title[Section~1.\insertsectionnumber \currentsection]{Chapter 1. Motivating Examples}
\title[\currentsection]{Chapter 4. Selecting and Checking Models}
\author[Model Selection]{MAST90084 Statistical Modelling Slides}
\vskip 1cm
\institute[Ch4]{Dennis Leung  \\ \vskip 0.3cm SCHOOL OF MATHEMATICS AND STATISTICS\\
\vskip 0.3cm
   THE UNIVERSITY OF MELBOURNE}
\vskip 1cm
\date[Dennis Leung \hspace{8pt}]{}


\begin{document}

\begin{frame}

Diagnostics techniques in R are interspersed in Faraway's book.

\end{frame}


\section{\S4.2 Diagnostics}
\begin{frame}{\S4.2 Diagnostics}
\begin{wideitemize}
\item
Diagnostic tools detect discrepancies between 
\begin{wideitemize}
\item the data and the fitted values;
\item a few data points and the rest
\end{wideitemize}

\item
Techniques  based on graphical presentations of
\begin{wideitemize}
\item  \textit{residuals}, (as in a usual residual plot against fitted values to spot trend or outliers)
\item  \textit{hat matrix},  (Provides some info about whether a fit greatly changes by omitting an outlying observation)
\item \textit{case deletion measures}.
\end{wideitemize}
\end{wideitemize}
\end{frame}

\subsection{\S4.2.1 Diagnostic tools for classical linear models}
\begin{frame}{\S4.2.1 Diagnostic tools for classical linear models}
Linear model (LM): 
$$\textbf{y}=X\mbox{\boldmath $\beta$}+\mbox{\boldmath $\varepsilon$}, \quad
\mbox{\boldmath $\varepsilon$}\sim N(\textbf{0},\sigma^2I)$$
where $X: n\times p$ matrix, and $\mbox{\boldmath $\beta$}: p\times 1$ vector parameter.
\begin{enumerate}
\item
\textbf{Least square (LS) estimator} of $\mbox{\boldmath $\beta$}$:  
$\hat{\mbox{\boldmath $\beta$}}=(X^TX)^{-1}X^T\textbf{y}$.

\item
\textbf{Hat matrix} $H= (h_{ij})_{1 \leq i , j \leq n} = X(X^TX)^{-1}X^T$. \vspace{8pt}

$\hat{\textbf{y}}=H\textbf{y}$, interpreted as the perpendicular \textit{projection} of $\textbf{y}$ onto the $p$-dimensional subspace
$\{X\mbox{\boldmath $\beta$}, \mbox{\boldmath $\beta$}\in\mathbb{R}^p\}$. \vspace{8pt}
 \item In many books $H$ is denoted as $P$ (as in ``projection"), but also denoted $H$ here since it transforms ${\bf y}$ into its ``hat" version $\hat{\bf y}$.
\end{enumerate}
\end{frame}

\begin{frame}{Hat matrix}
\begin{wideitemize}


\item A projection matrix: \textit{symmetric} and \textit{idempotent}. ($H^2=H$).


\item
Determine the \textbf{leverage} of $\bf y$  on the 
fitted values $\hat{\textbf{y}}$, {\bf via the $X$ data } .
\item
$h_{ij}$:  shows the amount of leverage or influence exerted on $\hat{y}_i$ by $y_j$.


\item The influence is due to $X$,  not $\textbf{y}$, b/c $H$ is constructed from $X$.

 

\end{wideitemize}
\end{frame}

\begin{frame}{Hat matrix, diagonals}
\begin{itemize}
\item
$h_{ii}$ reflects the influence of $y_i$ on $\hat{y}_i$, and
\begin{equation} \label{MD_rep}
h_i = n^{-1}+ (n-1)^{-1} MD_i,
\end{equation}
where $x_i$ is the $i$-th row of $X$, $S = (n-1)^{-1}X^TX$ and $MD_i$ is the Mahalanobis distance
\[
MD_i = (x_i - \bar{x})^TS^{-1}(x_i - \bar{x}).
\] 
Roughly, it measures how ``outlying" the i-th data point is, as far as covariates are concerned.  (Seber and Lee (S \& L), ``Linear Regression Analysis",  2003, p.269)




\end{itemize}

\end{frame}

\begin{frame}{Hat matrix, diagonals}

\begin{wideitemize}

\item Since
\[
h_{ii} = \sum_{j=1}^n h_{ij}^2 = h_{ii}^2  + \sum_{j \neq  i } h_{ij}^2,
\]
we have $h_{ii} \geq h_{ii}^2$, and hence $h_{ii}\leq 1$. It follows from \eqref{MD_rep} that
\[
n^{-1} \leq  h_{ii}\leq 1
\]
\item
Since $\displaystyle \mbox{rank}(H)=\sum_{i=1}^n h_{ii}=p$,  the average diagonal element is 
$\frac{p}{n}$.

\item
Rule of thumb: a variable $i$ for which $h_{ii}>\frac{2p}{n}$ holds is considered a \textbf{high-leverage
point}. 


\end{wideitemize}

\end{frame}

\begin{frame}{Residual vectors}
\begin{wideitemize}
\item
Residual vector:
\[\textbf{r}= (r_1,\cdots, r_n)^T = \textbf{y}-\hat{\textbf{y}} = (I - H){\bf y},\] where $\hat{\textbf{y}}=X\hat{\mbox{\boldmath $\beta$}}$
is the \textbf{fitted value} vector. It shows which data points are ill-fitting.

\item Also recall a consistent estimate of $\sigma^2$ is $ \hat{\sigma}^2=\frac{\textbf{r}^T
\textbf{r}}{n-p}$

\item$\text{Cov}({\bf r}) = \sigma^2 (I  - H)$,  since $I- H$ is also idempotent 
\[
 \Rightarrow var(r_i) = \sigma^2 (1 - h_{ii})
 \]

\item
\textbf{$i$-th standardized residual}: 
$\displaystyle r_i^*=\frac{r_i}{\sqrt{1-h_{ii}}}$ 
\item  \textbf{$i$-th studentized residual} (S \& L, p.267): $r_i^{(t)}=\frac{r_i^*}{\hat{\sigma}}=\frac{r_i}{\hat{\sigma}\sqrt{1-h_{ii}}}$
 
\end{wideitemize}
\end{frame}

\begin{frame}{Role of $H$ and ${\bf r}$ via the lens of omitting single observations}
\begin{wideitemize}

\item Let $\hat{\mbox{\boldmath $\beta$}}_{(i)}$ be the  LS estimate of $\mbox{\boldmath $\beta$}$ based on 
the data excluding the $i$-th observation

\item The quantity
\[
\Delta_i\hat{\boldsymbol \beta}=\hat{\boldsymbol \beta}-\hat{\mbox{\boldmath $\beta$}}_{(i)}
=\frac{(X^TX)^{-1}\textbf{x}_ir_i}{1-h_{ii}}
\]
gauges the effect of omitting observation $i$.  (Refer to S\& L, p.268 for a proof)

\item
$\Delta_i\hat{\mbox{\boldmath $\beta$}}$ increases with increasing residual $r_i$ and increasing $h_{ii}$. 

\item Obviously, $h_{ii}$ only carries effect of $x_i$ and $r_i$ should also carry the effect of $y_i$.

\end{wideitemize}

\end{frame}
%
%\begin{frame}
%
%Now generalize these to GLM
%\end{frame}

\subsection{\S4.2.2 Generalised hat matrix}
\begin{frame}{\S4.2.2 Generalised hat matrix}
\begin{wideitemize}

\item
Consider $Y$ having $k$ categories, $\mbox{\boldmath $\mu$}_i=\textbf{h}(Z_i\mbox{\boldmath $\beta$})$,
$i=1,\cdots,g$, and $\mbox{\boldmath $\mu$}_i$ is a $q\times 1$ vector, $q=k-1$.
\item
The IRLS estimate, at convergence, satisfies
\[
\hat{\mbox{\boldmath $\beta$}}=\underbrace{(Z^TW(\hat{\mbox{\boldmath $\beta$}})Z)^{-1}}_{F(\hat{\boldsymbol \beta})^{-1} }Z^TW(
\hat{\mbox{\boldmath $\beta$}})\tilde{\textbf{y}}(\hat{\mbox{\boldmath $\beta$}}) 
,
\]
where $\tilde{\textbf{y}}(\hat{\boldsymbol \beta})=Z\hat{\mbox{\boldmath $\beta$}}
+(D^{-1}(\hat{\mbox{\boldmath $\beta$}}))^T(\textbf{y}-\mbox{\boldmath $\mu$}(\hat{\mbox{\boldmath $\beta$}}))$.
\item
The \textbf{generalised hat matrix} is the $gq\times gq$ matrix
\[
H=(W^{1/2}(\hat{\mbox{\boldmath $\beta$}}))^T Z \cdot
(Z^TW(\hat{\mbox{\boldmath $\beta$}})Z)^{-1} \cdot
Z^TW^{1/2}(\hat{\mbox{\boldmath $\beta$}})
\]
\end{wideitemize}
\end{frame}

\begin{frame} {Generalized hat matrix }


\begin{wideitemize}
\item
$H$: also idempotent and symmetric, and 
$\mbox{tr}(H)=\mbox{rank}(H)$.
\item
When  $q=1$,
\begin{enumerate}
\item
$0\leq h_{ii}\leq 1$, high values of $h_{ii}$ (close to 1) correspond to extreme points in the design space.
\item
$H$ now depends on both $X$ and $\hat{\textbf{y}}$. Extreme points in the design matrix does not
necessarily have a high value of $h_{ii}$.
\end{enumerate}
\item 
When $q>1$,
\begin{enumerate}
\item
Consider the blocks $H_{ij}$ of $H=(H_{ij})$, $i,j=1,\cdots,g$, where
$H_{ij}$ is a $q\times q$ matrix.
\item
The sub-matrix $H_{ii}$ corresponds to the $i$-th observation. 
\item $\mbox{det}(H_{ii})$ or
$\mbox{tr}(H_{ii})$ may be used as an indicator for the leverage of $\textbf{y}_i$.
\end{enumerate}
\end{wideitemize}
\end{frame}


\begin{frame}{Example: vaso constriction}
\begin{wideitemize}
\item Finney (1947) is a study on the occurrence of a reflex vaso constriction (binary response) in the skin of the digits after taking a single deep breath. 

\item The binary response (vaso constriction)  is regressed on  volume of air (VOLUME) and inspiration rate (RATE), using a binary logit model. 


\end{wideitemize}

\end{frame}

\begin{frame}{Example: vaso constriction}
\begin{itemize}
\item Dataset available from the R package \texttt{robcbi}, or other \href{https://www.stat.auckland.ac.nz/~wild/data/Rdatasets/}{\textcolor{red}{repos}}:

\includegraphics[height=7cm]{constriction}
\end{itemize}

\end{frame}

\begin{frame}{vasoconstriction : ``H plot"}
\includegraphics[height=6cm]{h_plot}
\\
$h_{31, 31}$ has a larger value, but VOLUME and RATE aren't really that extreme compared to others as seen in the data table. 
\end{frame}

\subsection{\S4.2.3 Residuals}
\begin{frame}{\S4.2.3 Pearson Residuals }
Generalised residuals in GLM should have the properties similar to that of the standardised residuals for the LM.
\begin{itemize}
\item
\textbf{Pearson residual}$ (q\times 1$ vector)
$$\textbf{r}_i^P=\Sigma_i^{-1/2}(\hat{\mbox{\boldmath $\beta$}})(\textbf{y}_i-\hat{\mbox{\boldmath $\mu$}}_i)$$
\item Comes from  Pearson goodness-of-fit statistic : 
$$\chi^2=\sum_{i=1}^g \chi_p^2(\textbf{y}_i,\hat{\mbox{\boldmath $\mu$}}_i)=\sum_{i=1}^g
(\textbf{r}_i^P)^T\textbf{r}_i^P=\sum_{i=1}^g (\textbf{y}_i-\hat{\mbox{\boldmath $\mu$}}_i)^T
\Sigma_i^{-1}(\hat{\mbox{\boldmath $\beta$}})(\textbf{y}_i-\hat{\mbox{\boldmath $\mu$}}_i)$$
where
\begin{enumerate}
\item
 $\Sigma_i (\hat{\mbox{\boldmath $\beta$}})$ is the estimated covariance matrix of $\textbf{y}_i$,
\item $\hat{\mbox{\boldmath $\mu$}}_i=\mbox{\boldmath $\mu$}_i(\hat{\mbox{\boldmath $\beta$}})$.
\end{enumerate}
\end{itemize}
\end{frame}

\begin{frame}{\S4.2.3 Pearson (standardized) residuals}
\begin{itemize}
\item
Analogous to the LM case,  in large sample situations
$$\mbox{cov}(\textbf{y}_i-\hat{\mbox{\boldmath $\mu$}}_i)\approx \Sigma_i^{1/2}(\hat{\mbox{\boldmath $\beta$}})
(I-H_{ii})(\Sigma_i^{1/2}(\hat{\mbox{\boldmath $\beta$}}))^T\quad \Rightarrow \quad
\mbox{cov}(\textbf{r}_i^P)\approx I-H_{ii}$$
\item \textbf{studentized version of the Pearson residual} vector: 

 $$\displaystyle \textbf{r}_{i,s}^P=(I-H_{ii})^{-1/2}\textbf{r}_i^P=(I-H_{ii})^{-1/2}
\Sigma_i^{-1/2}(\hat{\mbox{\boldmath $\beta$}})(\textbf{y}_i-\hat{\mbox{\boldmath $\mu$}}_i)$$

\end{itemize}
\end{frame}


\begin{frame} {Deviance residual}

\begin{itemize}
\item
\textbf{Deviance residual} ($q=1$ is assumed): 
$$r_i^D=\mbox{sign}(y_i-\hat{\mu}_i)\sqrt{\chi_D^2(y_i, \hat{\mu}_i)}$$
Note $\chi_D^2(y_i, \hat{\mu}_i)=2\{\ell_i(y_i)-\ell_i(\hat{\mu}_i)\}$ and
$\displaystyle D=\sum_{i=1}^g \chi_D^2(y_i,\hat{\mu}_i)$ is the deviance of the model.


\end{itemize}

\end{frame}

%\begin{frame} {Anscombe residuals}
%
%
%Mainly to deal with the skewness of 
%Pearson residuals  when $n_i$ are small. Won't cover for now; see McCullagh and Nelder (1989) for a more detailed exposition. 
%
%
%\end{frame}






\begin{frame}{vasoconstriction : residual plot}
\includegraphics[height=7cm]{residualplot} \\
($i = 4, 18$ can be identified as outliers)
\end{frame}


\begin{frame}{vasoconstriction : QQ plot for $r^P_{i,s}$}
\includegraphics[height=6cm]{qqplot} \\
Standardized residuals $r_{i,s}^P$ should be somewhat approximately $N(0, 1)$ if we have a good fit; $ i = 4, 18$ are again far from the rest. 
\end{frame}

\subsection{\S4.2.4 Case deletion}
\begin{frame}{\S4.2.4 Case deletion (In linear model) }
\begin{itemize}
\item
The  difference $\Delta_i\hat{\mbox{\boldmath $\beta$}} = \hat{\mbox{\boldmath $\beta$}}-
\hat{\mbox{\boldmath $\beta$}}_{(i)}$ indicates the influence of the $i$-th observation $(\textbf{y}_i, Z_i)$ on the estimation of
$\mbox{\boldmath $\beta$}$.
\item
In particular, 
$$c_i=(\hat{\mbox{\boldmath $\beta$}}_{(i)}-\hat{\mbox{\boldmath $\beta$}})^T
\mbox{cov}(\hat{\mbox{\boldmath $\beta$}})^{-1}(\hat{\mbox{\boldmath $\beta$}}_{(i)}-
\hat{\mbox{\boldmath $\beta$}})$$ 
is a measure of such influence.
\item {\bf Cook distance}(1977): 
\[ \frac{c_i}{\dim (\hat{\mbox{\boldmath $\beta$}})}=\frac{1}{\hat{\sigma}^2
\cdot \dim (\hat{\mbox{\boldmath $\beta$}})}(\hat{\mbox{\boldmath $\beta$}}_{(i)}-\hat{\mbox{\boldmath $\beta$}})^T
(X^TX)(\hat{\mbox{\boldmath $\beta$}}_{(i)}-\hat{\mbox{\boldmath $\beta$}})^T\]

\end{itemize}
\end{frame}

\begin{frame}{In GLM}
\begin{itemize}


\item
In GLM, it can be computationally expensive to compute the leave-one-out estimates 
$\hat{\mbox{\boldmath $\beta$}}_{(i)}$, for every $i=1,\cdots, g$. 

\item
One-step approximation can be used:
$$\hat{\mbox{\boldmath $\beta$}}_{(i),1}=F_{(i)}^{-1}(\hat{\mbox{\boldmath $\beta$}})Z_{(i)}^T
W_{(i)}(\hat{\mbox{\boldmath $\beta$}})\tilde{\textbf{y}}(\hat{\mbox{\boldmath $\beta$}})$$
where reduced Fisher matrix $F_{(i)}(\hat{\mbox{\boldmath $\beta$}})=F(\hat{\mbox{\boldmath $\beta$}})
-Z_i^TW_i(\hat{\mbox{\boldmath $\beta$}})Z_i$ and
$$Z_{(i)}^TW_{(i)}(\hat{\mbox{\boldmath $\beta$}})\tilde{\textbf{y}}(\hat{\mbox{\boldmath $\beta$}})
=Z^TW(\hat{\mbox{\boldmath $\beta$}})\tilde{\textbf{y}}(\hat{\mbox{\boldmath $\beta$}})-
Z_i^TW_i(\hat{\mbox{\boldmath $\beta$}})\tilde{\textbf{y}}_i(\hat{\mbox{\boldmath $\beta$}})$$


\end{itemize}

\end{frame}


\begin{frame} {Generalized Cook's distance}
\begin{itemize}
\item In GLM,  one can then use an approximate \textbf{generalised Cook's measure} {\footnotesize
$$c_{i,1}=(\hat{\mbox{\boldmath $\beta$}}_{(i),1}-\hat{\mbox{\boldmath $\beta$}})^T
\mbox{cov}(\hat{\mbox{\boldmath $\beta$}})^{-1}(\hat{\mbox{\boldmath $\beta$}}_{(i),1}-
\hat{\mbox{\boldmath $\beta$}})=(\textbf{r}_i^P)^T(I-H_{ii})^{-1}H_{ii}(I-H_{ii})^{-1}\textbf{r}_i^P$$}
\hspace*{-4pt}to assess the influence of the $i$-th observation $(\textbf{y}_i, Z_i)$ on estimating 
${\mbox{\boldmath $\beta$}}$.


\item The second equality above an be easily derived using the alternative expression (Hennevogl \& Kranert, 1988):
$$\hat{\mbox{\boldmath $\beta$}}_{(i),1}=\hat{\mbox{\boldmath $\beta$}}-
F^{-1}(\hat{\mbox{\boldmath $\beta$}})Z_i^TW_i^{1/2}(\hat{\mbox{\boldmath $\beta$}})
(I-H_{ii})^{-1}\Sigma_i^{-1/2}(\hat{\mbox{\boldmath $\beta$}})(\textbf{y}_i-\mbox{\boldmath $\mu$}_i
(\hat{\mbox{\boldmath $\beta$}}))$$
%and the fact 
%\[
%\text{cov} (\hat{\boldsymbol \beta})  = 
%\]
\item $c_{i,1}$ here can be readily computed from previously defined diagnostic elements, without having to run any iterative procedures again.
\end{itemize}

\end{frame}

\begin{frame}{Cook distance for Vaso Constriction (Ex 4.6 in F\& T)}
\includegraphics[height=7cm]{cookdistance}
\end{frame}

\end{document}






%
%
%
%\begin{frame}{\large \S2.1.3 GLM for normal, Gamma \& inverse Gaussian responses (2)}
%\begin{itemize}
%\item
%When $y$ has an inverse Gaussian distribution $\mbox{IG}(\mu,\sigma^2)$ (useful for modelling life time etc.), the pdf is
%\begin{eqnarray*}
%f(y|\mu,\sigma^2) &\!\!\!=\!\!\!&\left(\frac{1}{2\pi\sigma^2 y^3}\right)^{1/2}\exp\left(-\frac{(y-\mu)^2}{2\sigma^2\mu^2y}
%\right), \quad\!\!\! y>0, \mu>0, \sigma^2>0 \\
%&\!\!\!=\!\!\!& \left(\frac{1}{2\pi\sigma^2 y^3}\right)^{1/2}\exp\left\{\frac{-y\frac{1}{2\mu^2}+\frac{1}{\mu}}{\sigma^2}
%-\frac{1}{2\sigma^2 y}\right\}
%\end{eqnarray*}
%\item
%Hence natural parameter $\theta=-\frac{1}{2\mu^2}$;  dispersion parameter $\phi=\sigma^2$.
%\begin{eqnarray*}
%b(\theta)=-\frac{1}{\mu}=-(-2\theta)^{1/2} \quad \Rightarrow \quad b^\prime(\theta)=(-2\theta)^{-1/2}=\mu \\
%\Rightarrow \; b^{\prime\prime}(\theta)\!=\!(-2\theta)^{-3/2}\!=\!\mu^3 \;\; \Rightarrow \;\;
%v(\mu)\!=\!\mu^3 \;\; \Rightarrow \;\; \mbox{Var}(y)\!=\!\sigma^2\mu^3
%\end{eqnarray*}
%The natural link is obtained by setting $\eta=\theta, \quad\Rightarrow\quad \eta=-\frac{1}{2\mu^2}$, i.e.
%$$\eta=\theta=(b^\prime)^{-1}(\mu)=-(2\mu^2)^{-1}.$$
%In practice $\eta=\mu^{-2}$ is used as the natural link.
%\end{itemize}
%\end{frame}
%
%\subsection{\S2.1.4 GLM for binary and binomial responses}\label{sec:2.1.4}
%\begin{frame}{\S2.1.4 GLM for binary and binomial responses (1)}
%\begin{itemize}
%\item
%Suppose $y_i|\textbf{x}_i \sim \mbox{Binomial}(m_i,\pi_i)$ and $\bar{y}_i=\frac{y_i}{m_i}$ are used as the
%response observations, $i=1,\cdots,n$. Then the pmf is {\footnotesize
%$$f(y_i|m_i, \pi_i,\textbf{x}_i)\!=\!{m_i\choose y_i}\pi_i^{y_i}(1-\pi_i)^{m_i\!-\!y_i}\!=\!{m_i\choose y_i}\exp\left\{
%m_i\bar{y}_i\log\frac{\pi_i}{1\!-\!\pi_i}+m_i\log (1\!-\!\pi_i)\right\}$$}
%\item
%Natural parameter $\theta_i=\log\frac{\pi_i}{1-\pi_i}$; dispersion $\phi=1$; weight $\omega_i=m_i$.
%\item
%Then $\mu_i^*=E(y_i|\textbf{x}_i)=m_i\pi_i=m_i\frac{\exp(\theta_i)}{1+\exp(\theta_i)}$
%{\small $$\Rightarrow \quad \mu_i=E(\bar{y}_i|\textbf{x}_i)=\pi_i=\frac{\exp(\theta_i)}{1+\exp(\theta_i)}$$}
%\item
%Also $b(\theta_i)=-\log (1-\pi_i)=\log (1+e^{\theta_i}) \quad \Rightarrow \quad 
%b^\prime(\theta_i)=\frac{e^{\theta_i}}{1+e^{\theta_i}}=\pi_i$
%{\small $$\Rightarrow\quad b^{\prime\prime}(\theta_i)=\frac{e^{\theta_i}}{(1+e^{\theta_i})^2}=\pi_i(1-\pi_i)
%=v(\pi_i)$$}
%$\Rightarrow\quad \mbox{Var}(\bar{y}_i|\textbf{x}_i)=\sigma^2(\mu_i)=\phi v(\mu_i)/\omega_i
%=m_i^{-1}\pi_i(1-\pi_i)$.
%\end{itemize}
%\end{frame}
%
%\begin{frame}{\S2.1.4 GLM for binary and binomial responses (2)}
%\begin{itemize}
%\item
%The natural link is $\eta_i=\theta_i=\log\frac{\pi_i}{1-\pi_i}=(b^\prime)^{-1}(\theta_i)$,
%called the \textbf{logit} (or \textbf{logistic}) \textbf{link}, resulting in the \textit{logistic model} 
%$\displaystyle \log\frac{\pi_i}{1-\pi_i}\!=\!\textbf{z}_i^T\!\mbox{\boldmath $\beta$}.$
%\item
%Other commonly used link functions in GLM with binomial responses:
%\begin{enumerate} 
%\item
%\textbf{Probit link} $\eta_i=\Phi^{-1}(\pi_i)$, where $\Phi(\cdot)$ is the N(0,1) cdf. This results in the
%\textit{probit model} $$\Phi^{-1}(\pi_i)=\textbf{z}_i^T\!\mbox{\boldmath $\beta$}.$$
%\item
%\textbf{Complementary log-log link} $\eta_i=g(\pi_i)=\log (-\log(1-\pi_i))$, resulting in the
%\textit{complementary log-log model} 
%$$\log (-\log(1-\pi_i))=\textbf{z}_i^T\!\mbox{\boldmath $\beta$}.$$
%Note the corresponding \textbf{response function} is $h(\eta_i)\!=\!1\!-\!\exp(\!-\!\exp(\!\eta_i\!)\!)$.
%\item
%\textbf{Complementary log link} $\eta_i=g(\pi_i)=-\log(1-\pi_i)$, resulting in the
%\textit{complementary log model}
%$$-\log(1-\pi_i)=\textbf{z}_i^T\!\mbox{\boldmath $\beta$}.$$
%\item
%Any inverse cdf can serve as a link function for GLM with binomial responses.
%\end{enumerate}
%\end{itemize}
%\end{frame}
%
%\begin{frame}{\S2.1.4 Over-dispersion in binomial GLM}
%\begin{itemize}
%\item
%Sometimes $\mbox{var}(\bar{y}_i|\textbf{x}_i) > m_i^{-1}\pi_i(1-\pi_i)$, so-called \textbf{over-dispersion},
%where $\bar{y}_i=m_i^{-1}y_i$ and $y_i$ is the number of ``successes" in $m_i$ Bernoulli trials.
%This implies the actual distribution of $y_i$ is not binomial. Two causes for this happening in practice.
%\begin{enumerate}
%\item
%Unobserved heterogeneity, i.e. different Bernoulli trials in the data have different probabilities of ``success".
%\item
%Positive correlation between individual Bernoulli trials.
%\end{enumerate}
%\item
%Since the actual distribution for $y_i$ is difficult to get in this situation, the simplest way to account
%for this over-dispersion is assuming
%$$\mbox{var}(\bar{y}_i|\textbf{x}_i) =\phi\cdot\frac{\pi_i(1-\pi_i)}{m_i},\quad
%\mbox{with $\phi>0$ an over-dispersion parameter.}$$
%\item
%\textbf{Remark}. Assuming a \textbf{beta-binomial distribution} for $y_i$ can be used to account for the
%over-dispersion.
%\end{itemize}
%\end{frame}
%
%\subsection{\S2.1.5 GLM for count responses}\label{sec:2.1.5}
%\begin{frame}{\S2.1.5 GLM for count responses}
%\begin{itemize}
%\item
%If $y_i|\textbf{x}_i \sim \mbox{Poisson}(\lambda=\mu_i)$ with $E(y_i|\textbf{x}_i)=\mu_i$
%and $\mbox{var}(y_i|\textbf{x}_i)=\mu_i$, it is easy to see the \textit{natural parameter} is
%$\theta_i=\log\mu_i$, and the \textit{natural link} is the \textbf{log-link} $\eta_i=\log\mu_i$, resulting in
%the \textbf{log linear model}:
%$$\log\mu_i=\eta_i=\textbf{z}_i^T\!\mbox{\boldmath $\beta$}, \quad \Rightarrow\quad
%\mu_i=\exp(\textbf{z}_i^T\!\mbox{\boldmath $\beta$})=\exp(\eta_i).$$ 
%\item
%A count response may not follow a Poisson distribution exactly, which can be accounted for
%by introducing an over-dispersion parameter $\phi$ such that $\mbox{var}(y_i|\textbf{x}_i)=\phi\mu_i$.
%\end{itemize}
%\end{frame}
%
%\section{\S2.2 Likelihood Inference}\label{sec:2.2}
%\begin{frame}{\S2.2 Likelihood Inference}
%\begin{itemize}
%\item
%Regression analysis with GLMs is based on \textbf{likelihood maximization}. 
%\item
%It includes parameter point estimation, confidence intervals, hypothesis testing, goodness-of-fit tests,
%and variable selection (or model selection).
%\item
%At the moment we assume the GLM under study is completely and correctly specified.
%This assumption can be relaxed later.
%\end{itemize}
%\end{frame}
%
%\subsection{\S2.2.1 Maximum likelihood estimation}\label{sec:2.2.1}
%\begin{frame}{\S2.2.1 Maximum likelihood estimation (1)}
%\begin{itemize}
%\item
%Consider the data with response $y_i$, covariates $\textbf{x}_i$ and design
%vector $\textbf{z}_i$, $i=1,2,\cdots, n$.
%\item
%Consider a GLM with $E(y_i|\textbf{x}_i)=\mu_i=h(\eta_i)=h(\textbf{z}_i^T\!\mbox{\boldmath $\beta$})$,
%where $h(\cdot)$ is the response function and $\mbox{\boldmath $\beta$}$ an unknown $p\times 1$ vector
%parameter.
%\item
%The focus is estimating $\mbox{\boldmath $\beta$}$ . The scale or dispersion parameter $\phi$ is
%assumed known or consistently estimated separately.
%\item
%Assume the design matrix $Z_{n\times p}=(\textbf{z}_1,\cdots,\textbf{z}_n)^T$ has (full) rank $p$.
%\item
%Now the log-likelihood function is {\footnotesize
%\begin{eqnarray*}
%\ell(\mbox{\boldmath $\beta$})&=&\sum_{i=1}^n\ell_i(\mbox{\boldmath $\beta$})
%=\sum_{i=1}^n f(y_i|\theta_i,\phi,\omega_i) \\
%&=&\sum_{i=1}^n \frac{y_i\theta_i-b(\theta_i)}{\phi}\cdot\omega_i
%+\mbox{terms not involving $\theta_i$'s} \\
%&=& \sum_{i=1}^n \frac{y_i\theta (\mu_i)-b(\theta(\mu_i))}{\phi}\cdot\omega_i
%=\sum_{i=1}^n \frac{y_i\theta (h(\textbf{z}_i^T\!\mbox{\boldmath $\beta$}))-
%b(\theta(h(\textbf{z}_i^T\!\mbox{\boldmath $\beta$})))}{\phi}\cdot\omega_i
%\end{eqnarray*}}
%\end{itemize}
%\end{frame}
%
%\begin{frame}{\S2.2.1 Maximum likelihood estimation (2)}
%\begin{itemize}
%\item
%The \textbf{score function} is {\footnotesize
%\begin{eqnarray*}
%\textbf{s}(\mbox{\boldmath $\beta$})&\!\!\!=\!\!\!& \frac{\partial \ell(\mbox{\boldmath $\beta$})}{\partial\mbox{\boldmath 
%$\beta$}}=\sum_{i=1}^n\frac{\partial \ell_i(\mbox{\boldmath $\beta$})}{\partial\mbox{\boldmath $\beta$}}
%=\sum_{i=1}^n  \frac{\partial \ell_i(\mbox{\boldmath $\beta$})}{\partial\theta_i}\cdot
%\frac{\partial \theta_i}{\partial\mu_i}\cdot \frac{\partial \mu_i}{\partial\eta_i} \cdot
%\frac{\partial \eta_i}{\partial\mbox{\boldmath $\beta$}}\\
%&\!\!\!=\!\!\!&\sum_{i=1}^n \left(\frac{y_i-b^\prime(\theta_i)}{\phi}\cdot\omega_i\right)\cdot
%\left[v(h(\textbf{z}_i^T\!\mbox{\boldmath $\beta$}))\right]^{-1}\cdot
% \frac{\partial h(\eta_i)}{\partial\eta_i}\cdot \textbf{z}_i \\
%&\!\!\!=\!\!\!& \sum_{i=1}^n \textbf{z}_i\cdot \frac{\partial h(\eta_i)}{\partial\eta_i}\cdot
%\left[\frac{\phi v(h(\textbf{z}_i^T\!\mbox{\boldmath $\beta$}))}{\omega_i}\right]^{-1}\cdot
%\left[y_i-\mu_i(\mbox{\boldmath $\beta$})\right] \\
%&\!\!\!\stackrel{denoted}{=} \!\!\!& \sum_{i=1}^n \textbf{z}_i D_i(\mbox{\boldmath $\beta$})
%\sigma_i^{-2}(\mbox{\boldmath $\beta$}) \left[y_i-\mu_i(\mbox{\boldmath $\beta$})\right] \\
%&\!\!\!=\!\!\!& \left[\textbf{z}_1, \cdots, \textbf{z}_n\right]  \left[\!\begin{array}{ccc}
%\!\!\!D_1(\mbox{\boldmath $\beta$})\!\! & & \\
%& \!\!\!\ddots\!\!\! & \\ & & \!\!\!D_n(\mbox{\boldmath $\beta$})\!\!\!\end{array}\!\right] 
%\left[\begin{array}{ccc}\!\!\!\!\sigma_1^2(\mbox{\boldmath $\beta$}) \!\!\! & & \\
%&\!\!\!\!\ddots \!\!\!\!& \\ & & \!\!\!\!\sigma_n^2(\mbox{\boldmath $\beta$})\!\!\!\end{array}\right]^{\!\!-1}
%\left[\begin{array}{c} \!\!\!y_1\!-\!\mu_1(\mbox{\boldmath $\beta$}) \!\!\!\! \\
%\!\!\!\vdots \!\!\!\! \\ \!\!\! y_n\!-\!\mu_n(\mbox{\boldmath $\beta$})\!\!\!\!\end{array}\!\right] \\
%&\!\!\!\stackrel{\mbox{matrix form}}{=} \!\!\!& Z^TD(\mbox{\boldmath $\beta$})\Sigma^{-1}(\mbox{\boldmath $\beta$})
%\left[\textbf{y}-\mbox{\boldmath $\mu$}(\mbox{\boldmath $\beta$})\right]
%\end{eqnarray*}}
%\end{itemize}
%\end{frame}
%
%\begin{frame}{\S2.2.1 Maximum likelihood estimation (3)}
%Supporting arguments for the derivation of $s(\mbox{\boldmath $\beta$})$
%\begin{itemize}
%\item
%$\displaystyle \mu_i=b^\prime(\theta_i) \quad \Rightarrow\quad \frac{\partial \mu_i}{\partial\theta_i}=
%b^{\prime\prime}(\theta_i)=v(\mu_i)=v(h(\textbf{z}_i^T\!\mbox{\boldmath $\beta$}))$
%\item
%Also $\displaystyle \frac{\partial \theta_i}{\partial\mu_i}=\left(\frac{\partial \mu_i}{\partial\theta_i}\right)^{-1}$ 
%and $\displaystyle \frac{\partial \eta_i}{\partial\mbox{\boldmath $\beta$}}=\textbf{z}_i$.
%\item
%{\footnotesize Denote $\displaystyle D_i(\mbox{\boldmath $\beta$})=\frac{\partial \mu_i}{\partial\eta_i}
%=\frac{\partial h(\eta_i)}{\partial\eta_i}\mid_{\eta_i=\textbf{z}_i^T\!\mbox{\boldmath $\beta$}}$,\;
%$\sigma_i^2(\mbox{\boldmath $\beta$})=\omega_i^{-1}\phi v(h(\textbf{z}_i^T\!\mbox{\boldmath $\beta$}))$
%and $w_i(\mbox{\boldmath $\beta$})=D_i^2(\mbox{\boldmath $\beta$})\sigma_i^{-2}(\mbox{\boldmath $\beta$})$,
%$i=1,\cdots, n$ (or $i=1,\cdots, g$ for grouped data).}
%\item
%{\footnotesize Denote $\displaystyle D(\mbox{\boldmath $\beta$})=\left[\begin{array}{ccc}
%\!\!\!D_1(\mbox{\boldmath $\beta$})\!\! & & \\
%& \!\!\!\ddots\!\!\! & \\ & & \!\!\!D_n(\mbox{\boldmath $\beta$})\!\!\!\end{array}\right]=\mbox{diag}
%\{D_1(\mbox{\boldmath $\beta$}), \cdots, D_n(\mbox{\boldmath $\beta$})\}$}
%\item
%{\footnotesize Denote $\displaystyle \Sigma(\mbox{\boldmath $\beta$})=\left[\begin{array}{ccc}
%\!\!\!\sigma_1^2(\mbox{\boldmath $\beta$})\!\! & & \\
%& \!\!\!\ddots\!\!\! & \\ & & \!\!\!\sigma_n^2(\mbox{\boldmath $\beta$})\!\!\!\end{array}\right]=\mbox{diag}
%\{\sigma_1^2(\mbox{\boldmath $\beta$}), \cdots, \sigma_n^2(\mbox{\boldmath $\beta$})\}$}
%\item
%{\footnotesize Denote $\displaystyle W(\mbox{\boldmath $\beta$})=\left[\begin{array}{ccc}
%\!\!\!w_1(\mbox{\boldmath $\beta$})\!\! & & \\
%& \!\!\!\ddots\!\!\! & \\ & & \!\!\!w_n(\mbox{\boldmath $\beta$})\!\!\!\end{array}\right]=\mbox{diag}
%\{w_1(\mbox{\boldmath $\beta$}), \cdots, w_n(\mbox{\boldmath $\beta$})\}$}
%\end{itemize}
%\end{frame}
%
%\begin{frame}{\S2.2.1 Maximum likelihood estimation (4)}
%\begin{itemize}
%\item
%The \textbf{Fisher information} matrix with respect to $\mbox{\boldmath $\beta$}$ is
%\begin{eqnarray*}
%F(\mbox{\boldmath $\beta$})&=& \mbox{cov}\left(\textbf{s}(\mbox{\boldmath $\beta$})\right)=
%\sum_{i=1}^n \textbf{z}_i D_i^2(\mbox{\boldmath $\beta$})\sigma_i^{-2}(\mbox{\boldmath $\beta$}) 
%\textbf{z}_i^T=\sum_{i=1}^n \textbf{z}_i w_i(\mbox{\boldmath $\beta$})\textbf{z}_i^T \\
%&=& \left[\textbf{z}_1, \cdots, \textbf{z}_n\right]  \left[\begin{array}{ccc}
%\!\!\!w_1(\mbox{\boldmath $\beta$})\!\! & & \\
%& \!\!\!\ddots\!\!\! & \\ & & \!\!\!w_n(\mbox{\boldmath $\beta$})\!\!\!\end{array}\right] 
%\left[\begin{array}{c} \textbf{z}_1^T \\ \vdots \\ \textbf{z}_n^T\end{array}\right]
%=Z^TW(\mbox{\boldmath $\beta$})Z
%\end{eqnarray*}
%\end{itemize}
%\end{frame}
%
%\begin{frame}{\S2.2.1 Maximum likelihood estimation (5)}
%\begin{itemize}
%\item
%The \textbf{observed information} matrix with respect to $\mbox{\boldmath $\beta$}$ is {\footnotesize
%\begin{eqnarray*}
%F_{\mbox{obs}}(\mbox{\boldmath $\beta$})& \!\!\!= \!\!\!& -\frac{\partial^2\ell(\mbox{\boldmath $\beta$})}{
%\partial \mbox{\boldmath $\beta$}\partial \mbox{\boldmath $\beta$}^T}=
%-\frac{\partial \textbf{s}(\mbox{\boldmath $\beta$})}{\partial \mbox{\boldmath $\beta$}^T} \\
%& \!\!\!= \!\!\!& -\sum_{i=1}^n \textbf{z}_i \frac{\partial \{D_i(\mbox{\boldmath $\beta$})\sigma_i^{-2}(
%\mbox{\boldmath $\beta$})\} }{\partial \mbox{\boldmath $\beta$}^T}[y_i-\mu_i(\mbox{\boldmath $\beta$})]
%+\sum_{i=1}^n \textbf{z}_i D_i(\mbox{\boldmath $\beta$})\sigma_i^{-2}(\mbox{\boldmath $\beta$}) \cdot
%\frac{\partial\mu_i}{\partial \eta_i}\cdot \frac{\partial\eta_i}{\partial \mbox{\boldmath $\beta$}^T} \\
%& \!\!\!= \!\!\!& -\sum_{i=1}^n \textbf{z}_i \frac{\partial}{\partial \mbox{\boldmath $\beta$}^T}
%\{D_i(\mbox{\boldmath $\beta$})\sigma_i^{-2}(\mbox{\boldmath $\beta$})\}\cdot
%[y_i-\mu_i(\mbox{\boldmath $\beta$})] +\sum_{i=1}^n \textbf{z}_i D_i^2(\mbox{\boldmath $\beta$})
%\sigma_i^{-2}(\mbox{\boldmath $\beta$}) \cdot \textbf{z}_i^T
%\end{eqnarray*}}
%\item
%Hence $E\left[F_{\mbox{obs}}(\mbox{\boldmath $\beta$})\right]=\sum_{i=1}^n \textbf{z}_i 
%D_i^2(\mbox{\boldmath $\beta$})\sigma_i^{-2}(\mbox{\boldmath $\beta$}) \textbf{z}_i^T
%=F(\mbox{\boldmath $\beta$})=\mbox{cov}\left(\textbf{s}(\mbox{\boldmath $\beta$})\right)$.
%\end{itemize}
%\end{frame}
%
%\begin{frame}{\S2.2.1 Maximum likelihood estimation (6)}
%\begin{itemize}
%\item
%When natural link is used, $\textbf{s}(\mbox{\boldmath $\beta$})$ and $F(\mbox{\boldmath $\beta$})$ can be
%simplified to
%\begin{eqnarray*}
%\textbf{s}(\mbox{\boldmath $\beta$})&= & \frac{1}{\phi}\sum_{i=1}^n \omega_i\textbf{z}_i
%[y_i-\mu_i(\mbox{\boldmath $\beta$})]=\frac{1}{\phi}Z^T\Omega \left[\textbf{y}-\mbox{\boldmath $\mu$}
%(\mbox{\boldmath $\beta$})\right] \\
%F(\mbox{\boldmath $\beta$}) &= & \frac{1}{\phi}\sum_{i=1}^n \omega_iv(\mu_i(\mbox{\boldmath $\beta$}))
%\textbf{z}_i\textbf{z}_i^T=\frac{1}{\phi}Z^T\Omega V(\mbox{\boldmath $\beta$})Z
%\end{eqnarray*}
%where $\Omega=\mbox{diag}\{\omega_1,\cdots,\omega_n\}$ and $V(\mbox{\boldmath $\beta$})
%=\mbox{diag}\{v(\mu_1),\cdots, v(\mu_n)\}$.
%\item
%For natural links, $\theta_i=\eta_i=\textbf{z}_i^T\mbox{\boldmath $\beta$} \quad \Rightarrow\quad
%\mu_i=b^\prime (\theta_i)=b^\prime(\textbf{z}_i^T\mbox{\boldmath $\beta$})\equiv 
%h(\textbf{z}_i^T\mbox{\boldmath $\beta$})$
%
%$\Rightarrow \quad D_i(\mbox{\boldmath $\beta$})=\frac{\partial h(\eta_i)}{\partial \eta_i}
%=b^{\prime\prime}(\textbf{z}_i^T\mbox{\boldmath $\beta$})=v(h(\textbf{z}_i^T\mbox{\boldmath $\beta$}))$
%
%$\Rightarrow \quad D_i(\mbox{\boldmath $\beta$})\sigma_i^{-2}(\mbox{\boldmath $\beta$})
%=v(h(\textbf{z}_i^T\mbox{\boldmath $\beta$}))\cdot\frac{\omega_i}{\phi
%v(h(\textbf{z}_i^T\mbox{\boldmath $\beta$}))}=\frac{\omega_i}{\phi}$ and
%
%$\Rightarrow \quad w_i(\mbox{\boldmath $\beta$})=D_i^2(\mbox{\boldmath $\beta$})
%\sigma_i^{-2}(\mbox{\boldmath $\beta$})=D_i(\mbox{\boldmath $\beta$})\cdot \frac{\omega_i}{\phi}
%=\frac{1}{\phi}\omega_i v(h(\textbf{z}_i^T\mbox{\boldmath $\beta$}))$.
%\end{itemize}
%\end{frame}
%
%\begin{frame}{\S2.2.1 Numerical computation of the MLE $\hat{\mbox{\boldmath $\beta$}}$ by IRLS}
%\begin{itemize}
%\item
%MLE $\hat{\mbox{\boldmath $\beta$}}$ is the one maximizing $\ell(\mbox{\boldmath $\beta$})$, and often 
%is obtained by solving the \textbf{estimating equation} $\textbf{s}(\mbox{\boldmath $\beta$})=0$, which uses 
%the \textbf{Fisher scoring method} for computation: {\scriptsize
%\begin{eqnarray*}
%\hat{\mbox{\boldmath $\beta$}}^{(k+1)}&\!\!\!\!\!\!=\!\!\!\!\!\! & \hat{\mbox{\boldmath $\beta$}}^{(k)}+
%F^{-1}(\hat{\mbox{\boldmath $\beta$}}^{(k)})\cdot \textbf{s}(\hat{\mbox{\boldmath $\beta$}}^{(k)}) \\
%&\!\!\!\!\!\!=\!\!\!\!\!\! & \hat{\mbox{\boldmath $\beta$}}^{(k)}+ 
%\left[Z^TW(\hat{\mbox{\boldmath $\beta$}}^{(k)})Z\right]^{-1}
% Z^TD(\hat{\mbox{\boldmath $\beta$}}^{(k)})\Sigma^{-1}(\hat{\mbox{\boldmath $\beta$}}^{(k)})
%\left[\textbf{y}-\mbox{\boldmath $\mu$}(\hat{\mbox{\boldmath $\beta$}}^{(k)})\right] \\
%&\!\!\!\!\!\!=\!\!\!\!\!\! & \hat{\mbox{\boldmath $\beta$}}^{(k)}+ 
%\left[Z^TW(\hat{\mbox{\boldmath $\beta$}}^{(k)})Z\right]^{-1}
% Z^TD^2(\hat{\mbox{\boldmath $\beta$}}^{(k)})\Sigma^{-1}(\hat{\mbox{\boldmath $\beta$}}^{(k)})
%D^{-1}(\hat{\mbox{\boldmath $\beta$}}^{(k)}) 
%\left[\textbf{y}-\mbox{\boldmath $\mu$}(\hat{\mbox{\boldmath $\beta$}}^{(k)})\right] \\
%&\!\!\!\!\!\!=\!\!\!\!\!\! & \left[Z^TW(\hat{\mbox{\boldmath $\beta$}}^{(k)})Z\right]^{-1}\left\{
%Z^TW(\hat{\mbox{\boldmath $\beta$}}^{(k)})Z \hat{\mbox{\boldmath $\beta$}}^{(k)}
%+Z^TW(\hat{\mbox{\boldmath $\beta$}}^{(k)}) D^{-1}(\hat{\mbox{\boldmath $\beta$}}^{(k)})
%\left[\textbf{y}-\mbox{\boldmath $\mu$}(\hat{\mbox{\boldmath $\beta$}}^{(k)})\right]\!\right\} \\
%&\!\!\!\!\!\!=\!\!\!\!\!\!& \left[Z^TW(\hat{\mbox{\boldmath $\beta$}}^{(k)})Z\right]^{-1}\left\{
%Z^TW(\hat{\mbox{\boldmath $\beta$}}^{(k)})\left[Z \hat{\mbox{\boldmath $\beta$}}^{(k)}
%+ D^{-1}(\hat{\mbox{\boldmath $\beta$}}^{(k)})
%\left[\textbf{y}-\mbox{\boldmath $\mu$}(\hat{\mbox{\boldmath $\beta$}}^{(k)})\right]\right]\right\} \\
%&\!\!\!\!\!\!\stackrel{denoted}{=}\!\!\!\!\!\! & \left[Z^TW(\hat{\mbox{\boldmath $\beta$}}^{(k)})Z\right]^{-1}
%Z^TW(\hat{\mbox{\boldmath $\beta$}}^{(k)})\cdot \tilde{\textbf{y}}^{(k)}
%\end{eqnarray*}}
%where $\tilde{\textbf{y}}^{(k)}=(\tilde{y}_1^{(k)},\cdots, \tilde{y}_n^{(k)})^T
%=Z \hat{\mbox{\boldmath $\beta$}}^{(k)}+ D^{-1}(\hat{\mbox{\boldmath $\beta$}}^{(k)})
%\left[\textbf{y}-\mbox{\boldmath $\mu$}(\hat{\mbox{\boldmath $\beta$}}^{(k)})\right]$
%is the \textbf{working observation vector} at iteration $k$, $k=0,1,2,\cdots$.
%\end{itemize}
%\end{frame}
%
%\begin{frame}{\S2.2.1 Implementing Fisher scoring for computing $\hat{\mbox{\boldmath $\beta$}}$}
%\begin{enumerate}
%\item
%Start with an initial estimate $\hat{\mbox{\boldmath $\beta$}}^{(0)}$, often chosen as the unweighted
%least squares (LS) estimate from $(g(y_i),\textbf{z}_i)$, $i=1,\cdots,n$, where $g(\cdot)$ is the link
%function (modified if $g(y_i)$ is undefined).
%\item
%For $k=0,1,2,\cdots$ compute 
%$$\hat{\mbox{\boldmath $\beta$}}^{(k+1)}=\left[Z^TW(\hat{\mbox{\boldmath $\beta$}}^{(k)})Z\right]^{-1}
%Z^TW(\hat{\mbox{\boldmath $\beta$}}^{(k)})\cdot \tilde{\textbf{y}}^{(k)}$$
%until convergence is achieved.
%\end{enumerate}
%Fisher scoring method in GLM estimation is also called the \textbf{iteratively reweighted least squares (IRLS) method}.
%\end{frame}
%
%\begin{frame}{\S2.2.1 Asymptotic properties associated with the MLE $\hat{\mbox{\boldmath $\beta$}}$}
%\begin{enumerate}
%\item
%\textit{Asymptotic existence and uniqueless}. The MLE $\hat{\mbox{\boldmath $\beta$}}$ exists
%and (locally) unique in probability 1.
%\item
%\textit{Consistency}. If $\mbox{\boldmath $\beta$}$ is the ``true'' value, then 
%$\hat{\mbox{\boldmath $\beta$}}\rightarrow \mbox{\boldmath $\beta$}$ in probability 1 (\textit{weak consistency}),
%or $\hat{\mbox{\boldmath $\beta$}}\rightarrow \mbox{\boldmath $\beta$}$ with probability 1 (or almost surely,
%\textit{strong consistency}).
%\item
%\textit{Asymptotic normality}. $\hat{\mbox{\boldmath $\beta$}}\stackrel{a}{\sim} \mbox{Normal}
%\left(\mbox{\boldmath $\beta$}, F^{-1}(\hat{\mbox{\boldmath $\beta$}})\right)$, equivalently
%$$F^{1/2}(\hat{\mbox{\boldmath $\beta$}})\left(\hat{\mbox{\boldmath $\beta$}}-
%\mbox{\boldmath $\beta$}\right)\stackrel{d}{\rightarrow} N(0,1) \quad\mbox{as $n\rightarrow\infty$}.$$
%This implies $\mbox{cov}(\hat{\mbox{\boldmath $\beta$}})\stackrel{a}{=}A(\hat{\mbox{\boldmath $\beta$}})
%= F^{-1}(\hat{\mbox{\boldmath $\beta$}})$. The proof uses the property $\textbf{s}(\mbox{\boldmath $\beta$})
%\stackrel{a}{\sim} \mbox{Normal}\left(0, F(\mbox{\boldmath $\beta$})\right)$.
%\item
%Estimating $\phi$: 
%$\displaystyle \hat{\phi}=\frac{1}{g-p}\sum_{i=1}^g\frac{(\bar{y}_i-\hat{\mu}_i)^2}{v(\hat{\mu}_i)/n_i}$
%is a consistent estimator of $\phi$. (This is true for grouped data only.)
%\end{enumerate}
%\end{frame}
%
%\subsection{\S2.2.2 Linear hypothesis testing in GLM}
%
%\begin{frame}{\S2.2.2 Testing linear hypotheses}
%\begin{itemize}
%\item
%Most testing problems concerning $\mbox{\boldmath $\beta$}$ have the form of \textit{linear hypothesis}:
%$$H_0: C\mbox{\boldmath $\beta$}=\mbox{\boldmath $\xi$} \quad\mbox{vs.}\quad
%H_1: C\mbox{\boldmath $\beta$}\neq\mbox{\boldmath $\xi$}$$
%where $C$ is a known matrix of full row rank $s\leq p$, and $\mbox{\boldmath $\xi$}$ is a known vector.
%\item
%A special case is 
%$$H_0: \mbox{\boldmath $\beta$}_r=0 \quad\mbox{vs.}\quad
%H_1: \mbox{\boldmath $\beta$}_r\neq 0$$
%where $\mbox{\boldmath $\beta$}_r$ is a sub-vector of $\mbox{\boldmath $\beta$}$. This is
%equivalent to comparing the full model with one of its sub-models.
%\end{itemize}
%\end{frame}
%
%\begin{frame}{\S2.2.2 Likelihood ratio (LR), Wald and score tests (1)}
%\begin{itemize}
%\item
%Three types of tests are available for testing $H_0: C\mbox{\boldmath $\beta$}=\mbox{\boldmath $\xi$}
%\;\mbox{vs.}\; H_1: C\mbox{\boldmath $\beta$}\neq\mbox{\boldmath $\xi$}$:
%
%\begin{center}
%\textbf{Likelihood ratio (LR) test}, \textbf{Wald test}, and \textbf{score test}.
%\end{center}
%\item
%A special case is 
%$$H_0: \mbox{\boldmath $\beta$}_r=0 \quad\mbox{vs.}\quad
%H_1: \mbox{\boldmath $\beta$}_r\neq 0$$
%where $\mbox{\boldmath $\beta$}_r$ is a sub-vector of $\mbox{\boldmath $\beta$}$. This is
%equivalent to comparing the full model with one of its sub-models.
%\end{itemize}
%\end{frame}
%
%\begin{frame}{\S2.2.2 Likelihood ratio (LR), Wald and score tests (2)}
%\begin{enumerate}
%\item
%\textbf{LR test statistic}: $\lambda=-2\left\{\ell(\tilde{\mbox{\boldmath $\beta$}}) -\ell(\hat{
%\mbox{\boldmath $\beta$}})\right\}$, where $\tilde{\mbox{\boldmath $\beta$}}$ is the MLE of
%$\mbox{\boldmath $\beta$}$ under $H_0$, and $\hat{\mbox{\boldmath $\beta$}}$ is the MLE of
%$\mbox{\boldmath $\beta$}$ under $H_1$; also $\hat{\mbox{\boldmath $\beta$}}$ is the unrestricted MLE of
%$\mbox{\boldmath $\beta$}$. $\lambda$ needs to be divided by $\hat{\phi}$ for over-dispersion models.
%\item
%\textbf{Wald test statistic}: $W=\left(C\hat{\mbox{\boldmath $\beta$}}-\mbox{\boldmath $\xi$}\right)^T
%\left[CF^{-1}(\hat{\mbox{\boldmath $\beta$}})C^T \right]^{-1}
%\left(C\hat{\mbox{\boldmath $\beta$}}-\mbox{\boldmath $\xi$}\right)$, where $\hat{\mbox{\boldmath $\beta$}}$
%is the unrestricted MLE of $\mbox{\boldmath $\beta$}$.
%\item
%\textbf{Score test statistic}: $U=\textbf{s}^T(\tilde{\mbox{\boldmath $\beta$}})F^{-1}(\tilde{\mbox{\boldmath $\beta$}})
%\textbf{s}(\tilde{\mbox{\boldmath $\beta$}})$, where $\tilde{\mbox{\boldmath $\beta$}}$
%is the MLE of $\mbox{\boldmath $\beta$}$ under $H_0$.
%\end{enumerate}
%\begin{itemize}
%\item
%Both Wald and score test statistics are properly defined for over-dispersion models.
%\item
%Under $H_0: C\mbox{\boldmath $\beta$}=\mbox{\boldmath $\xi$}$ with $\mbox{rank}(C)=s$, $\lambda, W$ and
%$U$ all follow a $\chi^2(s)$ distribution asymptotically under `regularity conditions'.
%\item
%$\lambda/\hat{\phi}$, $W$ and $U$ can be used for \textbf{analysis of deviance} (i.e. comparing nested
%models) as a special case of testing linear hypotheses.
%\end{itemize}
%\end{frame}
%
%\subsection{\S2.2.3 Deviance, residuals and goodness-of-fit measure in GLM} \label{sec:2.2.3}
%\begin{frame}{\S2.2.3 Model adequacy measure in GLM (1)}
%\begin{itemize}
%\item
%In linear regression, goodness-of-fit or adequacy of a model is measured by the sum of squared residuals (or errors)
%$\mbox{SSE}\!=\!\sum_{i=1}^n(y_i-\hat{\mu}_i)^2$. SSE is no longer a reasonable measure if $\mbox{Var}(y)$ 
%varies with $\mu$.
%\item
%A generalisation of SSE is the \textbf{Pearson $\chi^2$ statistic}, taking the form
%$$\chi^2=\sum_{i=1}^g \frac{(\bar{y}_i-\hat{\mu}_i)^2}{v(\hat{\mu}_i)/n_i}$$
%where $\bar{y}_i$ is the group mean, $v(\hat{\mu}_i)$ is the estimated variance function, and $n_i$ is the size of group $i$.
%\end{itemize}
%\end{frame}
%
%\begin{frame}{\S2.2.3 Model adequacy measure in GLM (2)}
%\begin{itemize}
%\item
%Pearson $\chi^2$ statistic is often an appropriate model adequacy measure (or discrepancy measure, or goodness-of-fit
%measure).
%\item
%Asymptotically (i.e. all $n_i$'s are large) $\chi^2/\phi$ follows a $\chi^2(g-p)$ distribution if the underlying GLM 
%provides adequate fit to the data. Here $p$ is the number of coefficient parameters in the GLM.
%\item
%If $n$ is large but $n_i$'s are small (in particular when $n_i=1$), it is dangerous to use the Pearson
%$\chi^2$ statistics to test model adequacy.
%\item
%If over-dispersion is present in the model, there is no need to test model adequacy.
%\end{itemize}
%\end{frame}
%
%\begin{frame}{\S2.2.3 Deviance measure in GLM (1)}
%\begin{itemize}
%\item
%Another model adequacy measure is based on the maximum log-likelihood for the data.
%\item
%It is twice the difference between the maximum log-likelihood achievable in a saturated model and that in
%the current model, i.e.
%$$D^*(\textbf{y}_n;\mbox{\boldmath $\hat{\mu}$})=2\left[\ln L(\mbox{\boldmath $\theta$}^*,\phi\mid\textbf{y}_n)
%-\ln L(\mbox{\boldmath $\hat{\theta}$},\phi\mid\textbf{y}_n)\right]$$
%which is called the \textbf{scaled deviance}. Here {\small \vspace{-3pt}
%\begin{equation*}
%L(\mbox{\boldmath $\theta$},\phi\mid\textbf{y}_n)=\exp\left\{\sum_{i=1}^n\frac{y_i\theta_i-b(\theta_i)}{\phi}\cdot
%\omega_i +\sum_{i=1}^nc(y_i,\phi, \omega_i)\right\} \vspace{-3pt}
%\label{eq-8.4}
%\end{equation*}}
%is the likelihood function for the data $\textbf{y}_n=(y_1,\cdots,y_n)^T$.
%\item
%$\mbox{\boldmath $\theta$}\!=\!(\theta_1,\cdots,\theta_n)^T$, $\mbox{\boldmath $\mu$}\!=\!(\mu_1,\cdots,\mu_n)^T
%\!=\!(b^\prime(\theta_1),\cdots, b^\prime(\theta_n))^T$, $\mbox{\boldmath $\hat{\theta}$}$ is the MLE of
%$\mbox{\boldmath $\theta$}$, and $\mbox{\boldmath $\theta$}^*=(\theta_1^*,\cdots, \theta_n^*)^T$ is such that 
%$\textbf{y}_n=(y_1,\cdots,y_n)^T
%=(b^\prime(\theta_1^*),\cdots,b^\prime(\theta_n^*))^T$.
%\end{itemize}
%\end{frame}
%
%\begin{frame}{\S2.2.3 Deviance measure in GLM (2)}
%\begin{itemize}
%\item
%We call $\mbox{\boldmath $\theta$}^*$ the MLE of $\mbox{\boldmath $\theta$}$ in the \textbf{saturated model}.
%\item
%In a saturated model, the number of independent parameters is the same as the number of independent observations.
%From the $L(\mbox{\boldmath $\theta$},\phi\mid\textbf{y}_n)$ expression we see that in this case, the MLE of 
%$\theta_i$ needs to satisfy $y_i=b^\prime(\theta_i)$.
%\item
%In GLM, $\theta_1,\cdots,\theta_n$ are determined by the regression parameters $\beta_1,\cdots,\beta_p$
%and the independent variables $x_1,\cdots,x_p$. So the MLE of $\beta_1,\cdots,\beta_p$ are to be obtained first, 
%the MLE $\mbox{\boldmath $\hat{\theta}$}$ are then obtained.
%\item
%The saturated model provides the best fit to the data because it has the same number of independent parameters
%and the observations.
%Maximum likelihood from the saturated model must be at least as large as that from a usual GLM (because $p<n$).
%Therefore, the scaled deviance $D^*(\textbf{y}_n;\mbox{\boldmath $\hat{\mu}$})\geq 0$ and measures
%the adequacy of fit of a GLM.
%\end{itemize}
%\end{frame}
%
%\begin{frame}{\S2.2.3 Deviance measure in GLM (3)}
%\begin{itemize}
%\item
%Consider a random sample $\textbf{y}_n=(y_1,\cdots,y_n)^T$ and a GLM as being formulated so far.
%We define 
%$$D(\textbf{y}_n;\mbox{\boldmath $\hat{\mu}$})=\phi D^*(\textbf{y}_n;\mbox{\boldmath $\hat{\mu}$})$$
%and call $D(\textbf{y}_n;\mbox{\boldmath $\hat{\mu}$})$ the \textbf{deviance} of the GLM under consideration.
%\item
%In most cases, $$\phi^{-1}D(\textbf{y}_n;\mbox{\boldmath $\hat{\mu}$})\stackrel{d}{\approx}\chi^2(n-p)$$ 
%asymptotically when the model for $Y$ is correct or adequate and `regularity conditions' are satisfied.
%\item
%But there are exceptions (some are important cases), e.g. when $n$ is large and all $n_i=1$.
%\item
%If over-dispersion is present, there is no need to test model adequacy by the deviance measure.
%\end{itemize}
%\end{frame}
%
%\begin{frame}{\S2.2.3 Deviances for several important distributions}
%\textbf{Deviances for several important distributions in GLM.}
%\begin{eqnarray*}
%\mbox{Normal}\!\!\! & N(\mu,\sigma^2): & D(\textbf{y}_n;\mbox{\boldmath $\hat{\mu}$})=
%\sum_{i=1}^n (y_i-\hat{\mu}_i)^2 \\
%\mbox{Poisson}\!\!\! &\mbox{Poi}(\mu): &  D(\textbf{y}_n;\mbox{\boldmath $\hat{\mu}$})=2\sum_{i=1}^n
%\left\{y_i\ln\frac{y_i}{\hat{\mu}_i}-(y_i-\hat{\mu}_i)\right\} \\
%\mbox{binomial}\!\!\! & \mbox{Bin}(m,\mu): & D(\textbf{y}_n;\mbox{\boldmath $\hat{\mu}$})=2\sum_{i=1}^n
%\left\{y_i\ln\frac{y_i}{\hat{\mu}_i}\!+\!(\!m_i\!-\!y_i\!)\ln\frac{m_i\!-\!y_i}{m_i\!-\!\hat{\mu}_i}\right\} \\
%\mbox{gamma} \!\!\!& \mbox{Gam}(\mu,\nu): & D(\textbf{y}_n;\mbox{\boldmath $\hat{\mu}$})=2\sum_{i=1}^n
%\left\{-\ln\frac{y_i}{\hat{\mu}_i}+\frac{y_i-\hat{\mu}_i}{\hat{\mu}_i}\right\}
%\end{eqnarray*}
%\end{frame}
%
%\begin{frame}{\S2.2.3 Analysis of deviance (1)}
%\begin{itemize}
%\item
%An analysis of variance (ANOVA) table can be constructed for each linear regression model as we have seen
%in a linear model course. 
%\item
%Similar to this, an \textbf{analysis of deviance (ANODEV)} table can be constructed for each GLM.
%\item
%Generally speaking, tests performed on ANOVA and ANODEV tables are likelihood ratio tests.
%\item
%However, there are some differences in interpreting ANOVA and ADODEV tables.
%\item
%Consider a $3\times 3$ factorial design, of two treatment factors A and B, with 2 replicates at each combination
%level of A and B. So there will be in total 18 observations for the response variable $y$.
%\end{itemize}
%\end{frame}
%
%\begin{frame}{\S2.2.3 Analysis of deviance (2)}
%If it is appropriate to fit the 18 $y$ observations by linear regression models, an ANOVA table 
%would typically look like the following:
%
%\begin{center} Analysis of Variance (ANOVA)\end{center}
%
%{\small
%\begin{tabular}{|cccccc|}\hline
%Model & df & SSE & $\Delta$SS & $\Delta$df & Effect Source \\ \hline
%1 & 17 & 1000 & - & -  & - \\ \hline
%A & 15 & 400 & $600=1000-400$ & $2=17-15$ & A \\ \hline
%A+B & 13 & 300 & $100=400-300$ & $2=15-13$ & B \\ \hline
%A+B+A:B & 9 & 100 & $200=300-100$ & $4=13-9$ & A:B \\ \hline 
%\end{tabular}}
%\begin{itemize}
%\item
%$\Delta$SS is the reduction of SSE as an additional term is added to the model.
%$\Delta\mbox{SS}\stackrel{d}{=}\chi^2(\Delta\mbox{df})$. Effects are assessed by
%$F$ tests.
%\item
%$\Delta$df is the reduction of the model degrees of freedom.
%\item
%The effects of A, B and A:B on $Y$ are orthogonal.
%\end{itemize}
%\end{frame}
%
%\begin{frame}{\S2.2.3 Analysis of deviance (3)}
%If it is more appropriate to fit the 18 $y$ observations by GLM, an ANODEV table would typically look like the following:
%
%\begin{center} Analysis of Deviance (ANODEV)\end{center}
%
%{\footnotesize
%\begin{tabular}{|cccccc|}\hline
%Model & df & Deviance & $\Delta D$ & $\Delta$df & Effect Source \\ \hline
%1 & 17 & 1200 & - & - & - \\ \hline
%A & 15 & 800 & $400\!=\!1200\!-\!800$ & $2\!=\!17\!-\!15$ & A ignoring B \\ \hline
%A+B & 13 & 500 & $300\!=\!800\!-\!500$ & $2\!=\!15\!-\!13$ & B eliminating A\\ \hline
%A+B+A:B & 9 & 200 & $300\!=\!500\!-\!200$ & $4\!=\!13\!-\!9$ & A:B eliminating \\ 
%& & & &  & A \& B \\ \hline 
%\end{tabular}}
%\begin{itemize}
%\item
%$\Delta D=D_0-D_1$ is the reduction of deviance between two nested models $M_0\subseteq M_1$.
%$\Delta D\stackrel{d}{\approx}\chi^2(\Delta\mbox{df})$ if $M_0$ is adequate.
%\item
%$\Delta$df is the reduction of the model degrees of freedom.
%\item
%Effects of A, B and A:B on $Y$ are not orthogonal, depend on their positions in the model, and
%are assessed by $\chi^2$ or $F$ tests.
%\end{itemize}
%\end{frame}
%
%%\section{\S8.4 Residuals}
%%\label{sec:residuals}
%
%\begin{frame}{\S2.2.3 GLM residuals (1)}
%\begin{itemize}
%\item
%In a linear regression model $y=E(y)+\varepsilon =\mu+\varepsilon=\sum_{j=1}^p x_j\beta_j+\varepsilon$,
%the residual is defined as $r_i=y_i-\hat{\mu}_i$ for an observation $y_i$, with $\hat{\mu}_i=\hat{y}_i$ being the 
%fitted value of $y_i$.
%The residual $r_i$ measures the goodness of fit to $y_i$ by the model if $\mbox{Var}(y)$ is constant.
%\item
%In GLM the residual of $y_i$ has to be defined differently to accommodate the cases
%where the variance function $v(\mu)$ is no longer constant.
%\item
%At least four types of residual are defined for GLM: \textbf{Pearson}, \textbf{deviance}, \textbf{working}, and
%\textbf{response}. All these residuals can be computed using the \texttt{R} function \texttt{residuals.glm()}.
%\end{itemize}
%\end{frame}
%
%\begin{frame}{\S2.2.3 GLM residuals (2)}
%\begin{itemize}
%\item
%\textbf{Pearson residual}: $r_{_{Pi}}=\frac{\bar{y}_i-\hat{\mu}_i}{\sqrt{v(\hat{\mu})/n_i}}$.
%
%Therefore $\sum_{i=1}^g r^2_{_{Pi}}=\sum_{i=1}^g \frac{(y_i-\hat{\mu}_i)^2}{v(\hat{\mu}_i)/n_i}=\chi^2$, the
%Pearson $\chi^2$ statistic. 
%\item
%\textbf{Deviance residual}: $r_{_{Di}}=\mbox{sign}(y_i-\hat{\mu}_i)\sqrt{d_i}$, \\ where $d_i$'s are such that
%$\displaystyle \sum_{i=1}^n d_i=D(\textbf{y}_n;\mbox{\boldmath $\hat{\mu}$})$. In particular, deviance residuals for
%the following distributions are:
%{\footnotesize
%\begin{eqnarray*}
%\mbox{Normal}\!\!\! & N(\mu,\sigma^2): & r_{_{Di}}\!=\!\mbox{sign}(y_i-\hat{\mu}_i)
%|y_i-\hat{\mu}_i|\!=\!y_i-\hat{\mu}_i \\
%\mbox{Poisson}\!\!\! &\mbox{Poi}(\mu): &  r_{_{Di}}\!=\!\mbox{sign}(y_i-\hat{\mu}_i)\sqrt{2\left\{y_i
%\ln\frac{y_i}{\hat{\mu}_i}-(y_i-\hat{\mu}_i)\right\}} \\
%\mbox{binomial}\!\!\! & \mbox{Bin}(m,\mu): & r_{_{Di}}\!=\!\mbox{sign}(y_i-\hat{\mu}_i)\sqrt{2
%\left\{y_i\ln\frac{y_i}{\hat{\mu}_i}\!+\!(\!m_i\!-\!y_i\!)\ln\frac{m_i\!-\!y_i}{m_i\!-\!\hat{\mu}_i}\right\}} \\
%\mbox{gamma} \!\!\!& \mbox{Gam}(\mu,\nu): & r_{_{Di}}\!=\!\mbox{sign}(y_i-\hat{\mu}_i)\sqrt{2
%\left\{-\ln\frac{y_i}{\hat{\mu}_i}+\frac{y_i-\hat{\mu}_i}{\hat{\mu}_i}\right\}}
%\end{eqnarray*}}
%\end{itemize}
%\end{frame}
%
%\begin{frame}{\S2.2.3 GLM residuals (3)}
%\begin{itemize}
%\item
%\textbf{Working residual}: $r_{_{Wi}}=(y_i-\hat{\mu}_i)\frac{d\eta_i}{d\mu_i}\mid_{\mu_i=\hat{\mu}_i}
%\approx g(y_i)-g(\hat{\mu}_i)$, \\
%where $g(\mu_i)$ is the link function and $\eta_i=\textbf{z}_i^T\mbox{\boldmath $\beta$}$ is the linear predictor in GLM.
%
%The working residual is used during computing the MLE of $\mbox{\boldmath $\beta$}$.
%\item
%\textbf{response residual}: $r_{_{Ri}}=y_i-\hat{\mu}_i$, \\
%which is not a right measure except when the GLM reduces to a linear model.
%\end{itemize}
%\end{frame}
%
%\begin{frame}[fragile] \frametitle{\S2.2.3 Example: Caesarian birth study (1)}
%For the data of Caesarian birth study we combine categories I and II of the response variable into
%one category \texttt{Infection}. Then fit logistic models to the reorganized data where the new response
%variable has only two categories \texttt{Infection} and \texttt{Non-infection}.
%
%\begin{scriptsize}
%\begin{Schunk}
%\begin{Sinput}
%> Caes.dat <- read.csv("D:/MAST90084/Example2-Caes.csv")
%> is.data.frame(Caes.dat)
%\end{Sinput} 
%\begin{Soutput}
%[1] TRUE
%\end{Soutput}
%\begin{Sinput}
%> Caes.dat
%\end{Sinput} 
%\begin{Soutput}
%  Infection total Antib RiskF Plan
%1         1    18     1     1    1
%2         0     2     1     0    1
%3        28    58     0     1    1
%4         8    40     0     0    1
%5        11    98     1     1    0
%6         0     0     1     0    0
%7        23    26     0     1    0
%8         0     9     0     0    0
%\end{Soutput}
%\end{Schunk}
%\end{scriptsize}
%\end{frame}
%
%\begin{frame}[fragile] \frametitle{\S2.2.3 Example: Caesarian birth study (2)}
%\begin{scriptsize}
%\begin{Schunk}
%\begin{Sinput}
%> Caes1=glm(Infection/total~Antib+RiskF+Plan,family=binomial, weight=total, data=Caes.dat)
%> summary(Caes1)
%\end{Sinput} 
%\begin{Soutput}
%Deviance Residuals: 
%       1         2         3         4         5         7         8  
% 0.26470  -0.15231  -0.78520   1.21563  -0.07162   1.49623  -2.56229  
%
%Coefficients:
%            Estimate Std. Error z value Pr(>|z|)    
%(Intercept)  -0.8207     0.4947  -1.659   0.0971 .  
%Antib        -3.2544     0.4813  -6.761 1.37e-11 ***
%RiskF         2.0299     0.4553   4.459 8.25e-06 ***
%Plan         -1.0720     0.4254  -2.520   0.0117 *  
%---
%Signif. codes:  0 '***' 0.001 '**' 0.01 '*' 0.05 '.' 0.1  ' ' 1
%
%(Dispersion parameter for binomial family taken to be 1)
%
%    Null deviance: 83.491  on 6  degrees of freedom
%Residual deviance: 10.997  on 3  degrees of freedom
%  (1 observation deleted due to missingness)
%AIC: 36.178
%Number of Fisher Scoring iterations: 5
%\end{Soutput}
%\end{Schunk}
%\end{scriptsize}
%\end{frame}
%
%\begin{frame}[fragile] \frametitle{\S2.2.3 Example: Caesarian birth study (3)}
%\begin{scriptsize}
%\begin{Schunk}
%\begin{Sinput}
%> anova(Caes1, test="Chi")
%\end{Sinput} 
%\begin{Soutput}
%Analysis of Deviance Table; Model: binomial, link: logit; Response: Infection/total
%Terms added sequentially (first to last)
%
%      Df Deviance Resid. Df Resid. Dev  Pr(>Chi)    
%NULL                      6     83.491              
%Antib  1   38.736         5     44.756 4.852e-10 ***
%RiskF  1   26.881         4     17.874 2.164e-07 ***
%Plan   1    6.878         3     10.997  0.008728 ** 
%\end{Soutput}
%\begin{Sinput}
%> resid(Caes1,type="deviance")
%\end{Sinput}
%\begin{Soutput}
%       1        2        3        4        5        7        8 
% 0.26470 -0.15231 -0.78520  1.21563 -0.07162  1.49623 -2.56229
%\end{Soutput}
%\begin{Sinput}
%> resid(Caes1,type="pearson")
%\end{Sinput}
%\begin{Soutput}
%       1        2        3        4        5        7        8 
% 0.27684 -0.10786 -0.78635  1.29465 -0.07141  1.38710 -1.99030
%\end{Soutput}
%\begin{Sinput}
%> sum(resid(Caes1,type="deviance")^2)    #[1] 10.99674
%> sum(resid(Caes1,type="pearson")^2)/(7-4)    #[1] 2.757724
%> 1-pchisq(sum(resid(Caes1,type="pearson")^2),3)   #[1] 0.04069089
%> 1-pchisq(sum(resid(Caes1,type="deviance")^2),3)   #[1] 0.01174351
%\end{Sinput}
%\end{Schunk}
%\end{scriptsize}
%\end{frame}
%
%\begin{frame}[fragile] \frametitle{\S2.2.3 Example: Caesarian birth study (4)}
%\begin{scriptsize}
%\begin{Schunk}
%\begin{Sinput}
%Caes2=glm(Infection/total~Antib+RiskF+Plan+RiskF:Plan,family=binomial, 
%                         weight=total, data=Caes.dat)
%summary(Caes2);  anova(Caes2,test="Chi")
%sum(resid(Caes2,type="pearson")^2)/(7-5)
%1-pchisq(sum(resid(Caes2,type="pearson")^2),2)
%1-pchisq(sum(resid(Caes2,type="deviance")^2),2)
%
%Caes3=glm(Infection/total~Antib+RiskF+Plan+RiskF:Plan,
%            family=binomial(link="probit"), weight=total, data=Caes.dat)
%summary(Caes3); anova(Caes3, test="Chi")
%
%Caes4=glm(Infection/total~Antib+RiskF+Plan+RiskF:Plan,
%           family=binomial(link="probit"), weight=total, data=Caes.dat[-6,])
%summary(Caes4); anova(Caes3, test="Chi")
%
%Caes5=glm(Infection/total~Antib+RiskF+Plan+RiskF:Plan,family=quasibinomial, 
%              weight=total, data=Caes.dat[-6,])
%summary(Caes5); anova(Caes5, test="F")
%
%sum(resid(Caes5,type="pearson")^2)/(7-5)
%
%Caes6=glm(Infection/total~Antib+RiskF+Plan+RiskF:Plan,family=quasibinomial, 
%           weight=total, data=Caes.dat, subset=c(1,2,3,4,5,7,8))
%summary(Caes6);  anova(Caes6, test="F")
%\end{Sinput} 
%\end{Schunk}
%\end{scriptsize}
%\end{frame}
%
%\section{\S2.3 Some extensions of GLM: quasi-likelihood and Bayesian method}
%\subsection{\S2.3.1 Quasi-likelihood models}
%\begin{frame}{\S2.3.1 Quasi-likelihood models (1)}
%\begin{itemize}
%\item
%Assuming that the response variable $y$ in GLM belongs to a specific exponential family
%confines the association between $E(y|\textbf{x})=\mu=h(\textbf{z}^T\mbox{\boldmath $\beta$})$
%and $\mbox{var}(y|\textbf{x})=\phi v(\mu)/\omega$ to be of a specific type, which sometimes is
%insufficient for practical use.
%\item
%Quasi-likelihood models drop the exponential family assumption and separate the mean and
%variance structure.
%
%No full distributional assumptions are necessary; only the first and second moments have
%to be specified. Under appropriate conditions, parameters can still be estimated consistently,
%and asymptotic inference is still possible.
%\end{itemize}
%\end{frame}
%
%\begin{frame}{\S2.3.1 Quasi-likelihood models (2)}
%\begin{itemize}
%\item
%Assume the mean and variance \textit{``working''} functions are
%\begin{eqnarray*}
%E(y_i|\textbf{x}_i)&=&\mu_i=h(\textbf{z}_i^T\mbox{\boldmath $\beta$}) \\
%\mbox{Var}(y_i|\textbf{x}_i) &=&\sigma^2(\mu_i)=\phi v(\mu_i)/\omega_i=\sigma^2_i(\mbox{\boldmath $\beta$})
%\end{eqnarray*}
%\item
%Also assume $\mu_i=h(\textbf{z}_i^T\mbox{\boldmath $\beta$})$ is correctly specified. 
%But $\sigma^2(\mu_i)=\phi v(\mu_i)/\omega_i$ is not necessarily correct.
%The true variance is assumed to be $\mbox{var}(y_i|\textbf{x}_i)=\sigma^2_{oi}(\textbf{x}_i)$.
%\item
%A \textbf{quasi-likelihood} (quasi log-likelihood more precisely) function can be constructed as
%$\displaystyle Q(\mbox{\boldmath $\mu$}; \textbf{y})=\sum_{i=1}^n\int_{y_i}^{\mu_i}\frac{y_i-t}{\omega_i^{-1}
%\phi v(t)}dt$.
%
%Note $-2Q(\mbox{\boldmath $\mu$}; \textbf{y})$ is called the \textbf{scaled quasi-deviance}.
%And the \textbf{quasi-deviance} is
%$$D(\textbf{y}, \mbox{\boldmath $\mu$})=-2\phi Q(\mbox{\boldmath $\mu$}; \textbf{y})
%=-2\sum_{i=1}^n\int_{y_i}^{\mu_i}\frac{y_i-t}{\omega_i^{-1} v(t)}dt$$
%\end{itemize}
%\end{frame}
%

\end{document}

